\section{Lattice tools and a prime-corank bound}\label{sec:typical}

The lattice-intersection argument of Section~\ref{sec:lattice} uses
three inputs: a uniform bound on corank modulo primes, a finite cover
of rational subspaces of low height, and a lattice decomposition that
removes directions of small determinant. We prove the prime-corank
bound for the biased bit model, adapting
\cite[Proposition~5.3]{SymmetricRate}, and recall the deterministic
lattice results with their precise hypotheses.

\subsection{Corank over all prime fields}

The principal-block argument of
\cite[Lemmas~2.2 and~2.3]{SymmetricRate} gives the following estimate over
any field in which the two bit values remain distinct. If \(A\) is a
symmetric \(d\times d\) matrix
whose upper-triangular entries are independent and have atoms at most
\(b\), and \(H\) is a fixed symmetric matrix over the same field, then
\begin{equation}\label{eq:typical-corank-tail}
 \Prob\bigl(\operatorname{corank}(A+H)\ge k\bigr)
 \le 2^d b^{k(k+1)/2},\qquad 1\le k\le d.
\end{equation}
Indeed, a symmetric matrix of rank \(r\) has a nonsingular principal
submatrix of size \(r\), over any field, including characteristic two.
One way to see this is to write it as \(Y C Y^{\mathsf T}\), where
\(Y\) has full column rank \(r\) and \(C\) is symmetric and
nonsingular, and then choose \(r\) independent rows of \(Y\).
For a fixed such principal block, reveal that block and all entries
joining it to its complement. If the full matrix has rank \(r\),
Schur complementation fixes all
\((d-r)(d-r+1)/2\) unrevealed entries. Their independence bounds the
conditional probability by \(b^{(d-r)(d-r+1)/2}\). Summing over all
principal sets of size at most \(d-k\) proves
Equation~\eqref{eq:typical-corank-tail}; the empty principal set handles
rank zero.

\begin{lemma}[Uniform prime corank]\label{lem:prime-corank}
Typically, every nonsingular bit matrix \(B\) of size \(m\in\mathcal I_N\)
satisfies
\[
 \operatorname{corank}_{\mathbb F_q}(B\bmod q)\le N^{3/5}
 \quad\text{for every prime }q.
\]
Consequently, its finite abelian group
\[
 G_B=\Z^m/B\Z^m
\]
can be generated by at most \(\lceil N^{3/5}\rceil\) elements.
The probability of nonsingularity together with failure is at most
\(2^{-\Omega_p(N^{6/5})}\).
\end{lemma}

\begin{proof}
For a fixed prime, Equation~\eqref{eq:typical-corank-tail} with
\(k=\lfloor N^{3/5}\rfloor+1\) gives the asserted probability bound.
If \(B\) is nonsingular over \(\R\), every prime with positive corank
divides the nonzero integer \(\det B\). Hadamard's inequality gives
\[
 1\le |\det B|\le m^{m/2}\le N^{N/2}.
\]
A union bound over all primes at most \(N^{N/2}\) therefore costs at
most \(2^{O(N\log N)}\), which is negligible compared with
\(2^{-\Omega_p(N^{6/5})}\). Finally,
\[
 \dim_{\mathbb F_q}(G_B/qG_B)
 =\operatorname{corank}_{\mathbb F_q}(B\bmod q).
\]
The structure theorem for finite abelian groups identifies the minimum
number of generators with the maximum of these dimensions.
\end{proof}

\subsection{Lattice height and a cover of rational subspaces}

A lattice is a discrete subgroup of a finite-dimensional Euclidean space,
considered in its real span. Its rank is the dimension of that span, and
its determinant is its Euclidean covolume. Thus, for a lattice basis
\(v_1,\ldots,v_r\),
\[
 \det L=\det\bigl(\langle v_i,v_j\rangle_{i,j=1}^r\bigr)^{1/2}.
\]
The rank-zero lattice has determinant one. A subspace
\(F\subseteq\R^m\) is \emph{rational} if it is spanned by integer
vectors; its \emph{height} is
\[
 H(F)=\det(F\cap\Z^m).
\]
In particular, \(H(0)=H(\R^m)=1\). Every such height is at least one:
the wedge of an integer lattice basis is a nonzero integer vector whose
Euclidean norm is the determinant.

We measure the difference between two subspaces by their \emph{rank distance}
\[
 d(U,V)=\dim U+\dim V-2\dim(U\cap V).
\]
This distance counts dimensions lost upon passing to their intersection;
it does not measure angles. The dimension formula gives
\(d(U^\perp,V^\perp)=d(U,V)\). The next theorem supplies a fixed finite
list approximating every subspace in the required height range.

\begin{theorem}[Low-height cover, {\cite[Theorem~4.1]{SymmetricRate}}]
\label{thm:height-cover}
Fix \(D,P>0\). For all sufficiently large \(N\) and every integer
\(0\le m\le N\), there is a list \(\mathcal C_{m,N}\) of rational
subspaces of \(\R^m\), depending only on \(m,N,D,P\), such that
\[
 |\mathcal C_{m,N}|\le 2^{N^2\ell^{-P}}.
\]
Every rational \(F\subseteq\R^m\) with
\(H(F)\le\exp(N\ell^D)\) satisfies
\[
 \min_{U\in\mathcal C_{m,N}}d(F,U)\le N\ell^{-P}.
\]
\end{theorem}

\subsection{Removing directions of small determinant}

The next decomposition separates a lattice from the directions in
which its determinant is small. A sublattice
\(L_0\subseteq L\) is \emph{saturated}, or \emph{primitive}, if
\(L_0=L\cap\Span L_0\). A basis of a saturated sublattice extends
to a basis of \(L\). If \(F=\Span L_0\), orthogonal projection onto
\(F^\perp\) therefore realizes the quotient as a lattice
\(L/L_0=\pi_{F^\perp}L\), and volume factorization gives
\begin{equation}\label{eq:orthogonal-quotient}
 \det L=\det L_0\,\det(\pi_{F^\perp}L).
\end{equation}
The same factorization holds for the full preimage in \(L\) of any
sublattice of the quotient. These facts hold in every rank, with the
rank-zero conventions already specified; see
\cite[Lemma~3.1]{SymmetricRate}.

Extremal sublattices and their canonical filtrations also occur in
Grayson's reduction theory of Euclidean lattices
\cite[Section~1]{Grayson}. We need only the following elementary
minimization statement.

\begin{lemma}[Determinant minimizer, {\cite[Lemma~3.6]{SymmetricRate}}]
\label{lem:minimizer}
For every lattice \(L\) and \(T>0\), the quantity
\[
 \frac{\det M}{T^{\rank M}},\qquad M\subseteq L\text{ a sublattice},
\]
attains its minimum at a saturated sublattice \(L_{\min}\). It satisfies
\[
 \det L_{\min}\le T^{\rank L_{\min}},\qquad
 \det M\ge T^{\rank M}
 \quad\text{for every sublattice }M\subseteq L/L_{\min}.
\]
\end{lemma}

\begin{proof}
The zero sublattice gives value one. For each positive rank \(r\), the
wedge of a basis belongs to the discrete lattice \(\bigwedge^r L\),
so there are only finitely many such wedge vectors of norm at most
\(T^r\). Thus the normalized determinants no greater than one have
only finitely many possible values, and the minimum is attained.
Saturation preserves rank and can only lower the determinant. Comparing
with the zero sublattice gives the first bound. For a sublattice \(M\)
of the quotient, its full preimage \(\widetilde M\subseteq L\) has
\[
 \det\widetilde M=\det L_{\min}\det M,
 \qquad
 \rank\widetilde M=\rank L_{\min}+\rank M.
\]
Minimality applied to \(\widetilde M\) gives the second bound.
\end{proof}

The point of removing \(L_{\min}\) is that lower bounds on all
sublattice determinants give an upper bound on the number of short
vectors in the remaining quotient.

\begin{lemma}[Projected point count]\label{lem:projected-count}
Let \(T>0\). Suppose a lattice \(\Lambda\) of rank \(j\ge1\) satisfies
\(\det M\ge T^{\rank M}\) for every sublattice \(M\subseteq\Lambda\).
Then, with \(s_0=T/(10(\log j+2))\),
\[
 \#\{x\in\Lambda:\lVert x\rVert_2\le R\}
 \le \frac32\exp(\pi R^2/s_0^2),\qquad R\ge0.
\]
In particular, for \(T=(\log N)^4\), \(j\le N\), and \(R=\sqrt N\),
this count is at most
\[
 2^{O(N/(\log N)^6)}=2^{o(N)}.
\]
The latter bound also holds in rank zero, where the count equals one.
\end{lemma}

\begin{proof}
For \(s>0\), write
\(\rho_s(\Lambda)=\sum_{x\in\Lambda}
\exp(-\pi\lVert x\rVert_2^2/s^2)\).
The Reverse Minkowski Theorem of Regev and Stephens-Davidowitz
\cite[Theorem~1.2]{ReverseMinkowski}, in the scaled form of
\cite[Theorem~3.4]{SymmetricRate}, gives
\(\rho_{s_0}(\Lambda)\le3/2\) under precisely the stated sublattice
hypotheses. Every vector of norm at most \(R\) contributes at least
\(\exp(-\pi R^2/s_0^2)\), yielding the displayed count; this is
\cite[Lemma~3.5]{SymmetricRate}. Finally,
\(s_0\ge c(\log N)^3\) for \(1\le j\le N\), which gives the claimed
exponent.
\end{proof}
